\chapter{Справочник с формули}
\label{app:formulas}

Този справочник събира на едно място формулите, които се използват най-често при
решаване на задачи. Всяко твърдение е доказано на съответното място в лекциите —
тук са само резултатите. Навсякъде $q=1-p$; за вероятностни модели с $p$ се
разбира $0\le p\le1$, а по-тесните области са отбелязани изрично.

\section{Комбинаторика}

\begin{center}
\small
\begin{tabular}{@{}llll@{}}
\toprule
& \textbf{редът има значение} & \textbf{редът няма значение} \\
\midrule
\textbf{без повторение} & $\dfrac{n!}{(n-k)!}$ (вариации) & $\dbinom{n}{k}$ (комбинации) \\[10pt]
\textbf{с повторение}   & $n^k$ & $\dbinom{n+k-1}{k}$ (стълбчета и звезди) \\
\bottomrule
\end{tabular}
\end{center}

\noindent
Тук избираме $k$ обекта измежду $n$, при $n,k\in\mathbb Z_{\ge0}$ и
$0\le k\le n$ за избор без повторение. За двата случая с повторение в
таблицата приемаме $n\ge1$ и $k\ge0$. Пермутациите на $n$ различни обекта са $n!$.

\begin{itemize}\itemsep2pt
  \item \emph{Мултиномен коефициент} (подреждания с повтарящи се букви): за
        $k_i\in\mathbb Z_{\ge0}$ и $\sum_i k_i=n$,
        $\dbinom{n}{k_1,\dots,k_m} = \dfrac{n!}{k_1!\,k_2!\cdots k_m!}$.
  \item \emph{Формула за включване и изключване:}
        $\displaystyle \Big|\bigcup_{i=1}^{n} A_i\Big|
          = \sum_i |A_i| - \sum_{i<j} |A_i \cap A_j| + \cdots + (-1)^{n-1}|A_1 \cap \dots \cap A_n|$.
  \item \emph{Брой сюрекции} от $k$-елементно върху $n$-елементно множество,
        $n\ge1$, $k\ge0$ цели:
        $\displaystyle \sum_{m=0}^{n} (-1)^m \binom{n}{m} (n-m)^k$.
        При $n=0$ има една сюрекция само ако $k=0$, и нито една при $k>0$.
  \item \emph{Свойства:} за $n\ge1$ и $0\le k\le n$, като в Pascal и
        absorption използваме стандартната нулева екстензия
        $\binom nk=0$ извън $0\le k\le n$,
        $\binom{n}{k} = \binom{n}{n-k}$;\quad
        $\binom{n}{k} = \binom{n-1}{k-1} + \binom{n-1}{k}$;\quad
        $k\binom{n}{k} = n\binom{n-1}{k-1}$;\quad
        $\sum_{k=0}^{n} \binom{n}{k} = 2^n$;\quad
        $\sum_{k=0}^{n}(-1)^k\binom{n}{k} = 0$ за $n \ge 1$.
  \item \emph{Бином на Нютон:} за $n\in\mathbb Z_{\ge0}$,
        $(a+b)^n = \sum_{k=0}^{n} \binom{n}{k} a^k b^{n-k}$.
        Извън $0\le k\le n$ приемаме $\binom nk=0$ само когато е обявена
        тази стандартна нулева екстензия.
\end{itemize}

\section{Основни вероятностни правила}

\begin{itemize}\itemsep2pt
  \item \emph{Класическа вероятност:} $\mathbb{P}(A) = \dfrac{|A|}{|\Omega|}$ при равновероятни изходи.
        \emph{Геометрична:} $\mathbb{P}(A) = \dfrac{|A|}{|\Omega|}$ с мярка (дължина, площ, обем).
  \item $\mathbb{P}(A^c) = 1 - \mathbb{P}(A)$;\quad
        $\mathbb{P}(A \cup B) = \mathbb{P}(A) + \mathbb{P}(B) - \mathbb{P}(A \cap B)$.
  \item \emph{Условна вероятност:} $\mathbb{P}(B \given A) = \dfrac{\mathbb{P}(A \cap B)}{\mathbb{P}(A)}$, при $\mathbb{P}(A) > 0$.
  \item \emph{Верижно правило} (когато всяко условящо сечение в знаменател е с
        положителна вероятност):
        $\mathbb{P}(A_1 \cap \dots \cap A_n) = \mathbb{P}(A_1)\mathbb{P}(A_2 \given A_1)\cdots
         \mathbb{P}(A_n \given A_1 \cap \dots \cap A_{n-1})$.
  \item \emph{Условна независимост:} при $\mathbb P(C)>0$,
        $\mathbb{P}(A \cap B \given C) = \mathbb{P}(A \given C)\mathbb{P}(B \given C)$
        е \emph{друго} свойство от независимостта на $A$ и $B$ — нито едното не влече другото.
        В примера с двата зара $A \perp\!\!\!\perp B$, но при условие $C$ е в сила $A \iff B$.
  \item \emph{Пълна вероятност по непрекъснат параметър:}
        $\mathbb{P}(A) = \int \mathbb{P}(A \given \Theta = \theta) f_\Theta(\theta)\,d\theta$ —
        така се смятат смеси (случаен параметър). Точковото условие означава избрана
        регулярна условна версия и е определено само за $\mathbb P_\Theta$-почти всяко
        $\theta$; стойностите върху нулево множество не влияят на интеграла. Вж.
        \crosslecture{supp:cond-indep}{допълнението в лекция 3}.
  \item \emph{Формула за пълната вероятност:} при пълна група $\{H_i\}$ с $\mathbb{P}(H_i)>0$,
        $\displaystyle \mathbb{P}(A) = \sum_i \mathbb{P}(A \given H_i)\,\mathbb{P}(H_i)$.
  \item \emph{Формула на Бейс:} при положителен знаменател и $\mathbb P(H_j)>0$,
        $\displaystyle \mathbb{P}(H_j \given A) = \frac{\mathbb{P}(A \given H_j)\mathbb{P}(H_j)}{\sum_i \mathbb{P}(A \given H_i)\mathbb{P}(H_i)}$.
  \item \emph{Независимост:} $A \perp\!\!\!\perp B \iff \mathbb{P}(A \cap B) = \mathbb{P}(A)\mathbb{P}(B)$.
  \item \emph{Две по две $\ne$ съвкупно:} събития може да са независими по двойки, без да са
        независими в съвкупност. Изисква се $\mathbb{P}(A_{i_1}\cap\dots\cap A_{i_k}) =
        \prod \mathbb{P}(A_{i_j})$ за \emph{всяко} подмножество, не само за двойките.
        Класически пример: два симетрични зара, $A=\{$първият е четен$\}$,
        $B=\{$вторият е четен$\}$, $C=\{$сборът е четен$\}$ — трите са две по две независими,
        но $C$ се определя от $A$ и $B$.
        В съвкупност: равенството важи за \emph{всяко} подмножество от индекси — това е
        по-силно от двойката по двойка.
\end{itemize}

\section{Дискретни разпределения}

\begin{center}
\footnotesize
\setlength{\tabcolsep}{4pt}
\begin{tabular}{@{}llccl@{}}
\toprule
\textbf{Разпределение} & $\mathbb{P}(X=k)$ & $\E X$ & $\Var X$ & $g_X(s)$ \\
\midrule
$\Ber(p)$ & $p^k q^{1-k}$, $k \in \{0,1\}$, $0\le p\le1$ & $p$ & $pq$ & $q+ps$ \\[3pt]
$\Bin(n,p)$ & $\dbinom{n}{k} p^k q^{n-k}$, $0 \le k \le n$, $n\in\mathbb Z_{\ge0}$ & $np$ & $npq$ & $(q+ps)^n$ \\[5pt]
$\Geo(p)$ & $q^k p$, $k \ge 0$, $0<p<1$ & $\dfrac{q}{p}$ & $\dfrac{q}{p^2}$ & $\dfrac{p}{1-qs}$ \\[7pt]
$\NegBin(r,p)$ & $\dbinom{r+k-1}{k} p^r q^k$, $k \ge 0$, $r\in\mathbb Z_{>0},\ 0<p<1$ & $\dfrac{rq}{p}$ & $\dfrac{rq}{p^2}$ & $\left(\dfrac{p}{1-qs}\right)^{\! r}$ \\[7pt]
$\Poi(\lambda)$ & $e^{-\lambda}\dfrac{\lambda^k}{k!}$, $k \ge 0$, $\lambda\ge0$ & $\lambda$ & $\lambda$ & $e^{\lambda(s-1)}$ \\[7pt]
$\Hyp(N,M,n)$ & $\dfrac{\binom{M}{k}\binom{N-M}{n-k}}{\binom{N}{n}}$ ${}^{(*)}$ & $n\dfrac{M}{N}$ & $n\dfrac{M}{N}\dfrac{N-M}{N}\dfrac{N-n}{N-1}$ & --- \\[7pt]
равномерно в $\{1,\dots,n\}$, $n\ge1$ & $\dfrac{1}{n}$ & $\dfrac{n+1}{2}$ & $\dfrac{n^2-1}{12}$ & $\dfrac{s(1-s^n)}{n(1-s)}$ \\
\bottomrule
\end{tabular}
\end{center}

\noindent
${}^{(*)}$ За $N\in\mathbb Z_{>0}$, $0\le M\le N$ и $0\le n\le N$; при $N>1$
е дадена дисперсията (при $N=1$ тя е $0$). Носителят на $\Hyp$ не е целият интервал $0..n$: изтеглените не могат да
надхвърлят наличните, нито да са твърде малко, тоест
$\max(0,\, n-N+M) \le k \le \min(M,\, n)$.

\noindent
$\Geo$ и $\NegBin$ в таблицата броят \emph{неуспехите} преди първия (съответно $r$-тия)
успех; $\NegBin(1,p) = \Geo(p)$. Множителят $\frac{N-n}{N-1}$ при $\Hyp$ е поправката за
теглене без връщане: при $N \to \infty$ той клони към $1$ и $\Hyp \to \Bin$.
При $p=0$ или $p=1$ биномното е дегенерирано съответно в $0$ или $n$;
$\Poi(0)$ е дегенерирано в $0$. За $\NegBin$ краят $p=1$ дава нула
неуспехи, а при $p=0$ времето до $r$-тия успех не е крайно.

\subsection*{Двете конвенции за $\Geo(p)$: броим ли последния опит}

Геометричното разпределение се среща в два варианта според това дали броим и
\emph{последния} (успешния) опит. Двата се различават само с изместване с единица, но
$\E$ и пораждащата функция изглеждат различно — оттам идват повечето грешки.

В таблицата по-долу $0<p<1$. При $p=1$ съответните величини са дегенерирани
($X\equiv0$, $Y\equiv1$), а при $p=0$ времето до успех не е крайно; затова
формулите с $1/p$ и последващите условни отношения не се пренасят към тези краища.

\begin{center}
\footnotesize
\setlength{\tabcolsep}{5pt}
\begin{tabular}{@{}lll@{}}
\toprule
 & $X$: брой \emph{неуспехи} & $Y = X+1$: брой \emph{опити} \\
 & (\emph{не} броим успешния) & (броим и успешния) \\
\midrule
носител & $k = 0,1,2,\dots$ & $k = 1,2,3,\dots$ \\[2pt]
$\mathbb{P}(\,\cdot\, = k)$ & $q^k p$ & $q^{k-1} p$ \\[2pt]
опашка & $\mathbb{P}(X \ge k) = q^{k}$ & $\mathbb{P}(Y > k) = q^{k}$ \\[2pt]
$\E$ & $\dfrac{q}{p}$ & $\dfrac{1}{p}$ \\[6pt]
$\Var$ & $\dfrac{q}{p^2}$ & $\dfrac{q}{p^2}$ (същата) \\[6pt]
$g(s)$ & $\dfrac{p}{1-qs}$ & $\dfrac{ps}{1-qs}$ \\[6pt]
липса на памет & $\mathbb{P}(X \ge m+k \given X \ge m) = \mathbb{P}(X \ge k)$
               & $\mathbb{P}(Y > m+k \given Y > m) = \mathbb{P}(Y > k)$ \\
\bottomrule
\end{tabular}
\end{center}

\begin{itemize}\itemsep2pt
  \item \emph{Преводът е едно изместване:} $Y = X + 1$, тоест $\E Y = \E X + 1$,
        $\Var Y = \Var X$ (изместването не мени дисперсията) и $g_Y(s) = s\,g_X(s)$.
        Обратно: $X = Y - 1$.
  \item \emph{Коя конвенция иска задачата:} ако въпросът е „колко пъти хвърляме зара
        \emph{до} първата шестица“, се брои и успешното хвърляне — това е $Y$ и
        $\E Y = 1/p = 6$. Ако въпросът е „колко хвърляния са неуспешни“ / „колко пъти
        пропускаме, преди да улучим“, това е $X$ и $\E X = q/p = 5$. Разликата е точно
        единица, така че отговорите $6$ и $5$ са двете конвенции, а не грешка в сметката.
  \item \emph{Бърза проверка:} при $p = 1$ успехът идва веднага, тоест $Y \equiv 1$, а
        $X \equiv 0$. Ако формулата дава $\E = 1/p$, броят се опити; ако дава $q/p$ —
        неуспехи.
  \item \emph{Внимание при липсата на памет:} за $Y$ неравенствата са \emph{строги}.
        С $\ge$ равенството се чупи: $\mathbb{P}(Y \ge m+k \given Y \ge m) = q^{k}$, докато
        $\mathbb{P}(Y \ge k) = q^{k-1}$.
  \item \emph{Същото при $\NegBin$:} таблицата дава броя неуспехи преди $r$-тия успех, а
        броят \emph{опити} до $r$-тия успех е $r + \NegBin(r,p)$, с очакване $r/p$.
  \item В тези записки и в \crosslecture{sec:geometric}{лекция 6} $\Geo(p)$ означава варианта $X$
        (неуспехи). Затова в задачата за колекционера по-долу се пише $T_i = 1 + \Geo(p_i)$.
\end{itemize}

\begin{itemize}\itemsep2pt
  \item \emph{Липса на памет} (само $\Geo$, при $0<p<1$ и $m,k\ge0$):
        $\mathbb{P}(X \ge m+k \given X \ge m) = \mathbb{P}(X \ge k)$.
  \item \emph{Мода:} за $0<p<1$, ако $(n+1)p$ не е цяло, единствената мода на $\Bin(n,p)$ е
        $\lfloor p(n+1)\rfloor$; ако $(n+1)p=m\in\mathbb Z$, модите са $m-1,m$.
        При $p=0$ модата е $0$, а при $p=1$ е $n$. За $\Poi(\lambda)$,
        $\lambda>0$, мода е $\lfloor\lambda\rfloor$ при нецяло $\lambda$, а при
        цяло $\lambda$ са $\lambda-1,\lambda$; при $\lambda=0$ модата е $0$.
  \item \emph{Сума със случаен брой събираеми:} при целочислено $N\ge0$, $N$ независимо
        от независимите еднакво разпределени $X_j$, и $S_N = \sum_{j=1}^N X_j$,
        пораждащите функции се \emph{съставят}. Ако $\E N,\E X_1<\infty$, важи
        формулата за средното; за дисперсията изискваме
        $\E N^2,\E X_1^2<\infty$:
        \[ g_{S_N}(s) = g_N\big(g_{X_1}(s)\big), \qquad
           \E S_N = \E N\,\E X_1, \qquad
           \Var S_N = \E N \Var X_1 + \Var N (\E X_1)^2 . \]
        Оттук прореждането на $\Poi(\lambda)$ с $\Ber(p)$ пада на един ред:
        $e^{\lambda(q+ps-1)} = e^{\lambda p(s-1)}$. Вж.
        \crosslecture{supp:random-sum}{допълнението в лекция 8}.
  \item \emph{Приближение на Поасон:} при голямо $n$ и малко $p$, $\Bin(n,p) \approx \Poi(np)$.
  \item \emph{Прореждане на Поасоново:} ако $N \sim \Poi(\lambda)$ и всеки от $N$-те обекта
        независимо се маркира с вероятност $p$, то броят маркирани $M \sim \Poi(\lambda p)$,
        броят немаркирани $K \sim \Poi(\lambda q)$, и $M \perp\!\!\!\perp K$. Вж.
        \crosslecture{supp:poisson-thinning}{допълнението в лекция 7}. При \emph{фиксиран} брой
        опити $n$ същото не важи: там $K = n - M$ е напълно определено от $M$.
  \item \emph{Сума от геометрични с различни параметри} (задачата за колекционера): ако
        $n\ge1$, $T_i$ е броят \emph{опити} на $i$-тия етап и $p_i = \frac{n-i+1}{n}$,
        то
        \[ \E T = n H_n = n\sum_{k=1}^{n}\frac1k \approx n\ln n + \gamma n, \qquad
           \Var T = n^2\sum_{k=1}^{n}\frac{1}{k^2} - n H_n
           \sim \frac{\pi^2}{6}n^2 . \]
        Схемата — етапи, линейност за очакването, независимост за дисперсията — е описана
        в \crosslecture{supp:collector}{допълнението в лекция 6}. Внимавайте: $\Geo$ брои
        неуспехи, а тук се броят опити, тоест $T_i = 1 + \Geo(p_i)$ и $\E T_i = 1/p_i$.
\end{itemize}

\begin{itemize}\itemsep2pt
  \item \emph{Бета интеграл} (полезен при осредняване по неизвестна вероятност):
        \[ \int_0^1 p^{k}(1-p)^{n-k}\,dp = \frac{k!\,(n-k)!}{(n+1)!} = \frac{1}{(n+1)\binom{n}{k}} . \]
        Следствие: ако $p \sim U(0,1)$ и при дадено $p$ броят успехи е $S_n \sim \Bin(n,p)$, то
        $\mathbb{P}(S_n = k) = \frac{1}{n+1}$ за всяко $k$ — $S_n$ е \emph{равномерно} върху
        $\{0,1,\dots,n\}$. Оттук и правилото за наследяване на Лаплас:
        \[ \mathbb{P}(\text{следващото е успех} \given S_n = k) = \frac{k+1}{n+2} . \]
\end{itemize}

\begin{itemize}\itemsep2pt
  \item \emph{Линейна рекурсия с неподвижна точка:} ако $p_{n+1} = a\,p_n + b$ с $|a| < 1$, то
        неподвижната точка е $p^{*} = \dfrac{b}{1-a}$ и
        \[ p_n = p^{*} + (p_0 - p^{*})\,a^{n} \;\xrightarrow[n\to\infty]{}\; p^{*} . \]
        Типично приложение: вероятност за връщане в изходно състояние при разходка по краен
        граф. Например при разходка по върховете на равностранен триъгълник
        $p_{n+1} = \frac{1-p_n}{2}$, тоест $a = -\frac12$, $b = \frac12$, $p^{*} = \frac13$ и
        $p_n = \frac13 + \frac23\left(-\frac12\right)^n \to \frac13$.
  \item \emph{Афинна замяна при нормално приближение:} ако $S \approx N(\mu, \sigma^2)$, то
        $aS + b \approx N(a\mu + b,\ a^2\sigma^2)$. Затова „брой загубени“ $= n - S$ и „брой думи“
        $= (n-S)/6$ се смятат направо, без нова ЦГТ.
\end{itemize}

\section{Непрекъснати разпределения}

\begin{center}
\footnotesize
\setlength{\tabcolsep}{4pt}
\begin{tabular}{@{}llccl@{}}
\toprule
\textbf{Разпределение} & $f_X(x)$ & $\E X$ & $\Var X$ & $M_X(t) = \E e^{tX}$ \\
\midrule
$U(a,b)$ & $\dfrac{1}{b-a}$, $a<x<b$, $a<b$ & $\dfrac{a+b}{2}$ & $\dfrac{(b-a)^2}{12}$ & $\dfrac{e^{tb}-e^{ta}}{t(b-a)}$ \\[7pt]
$\Exponential(\lambda)$ & $\lambda e^{-\lambda x}$, $x>0$, $\lambda>0$ & $\dfrac{1}{\lambda}$ & $\dfrac{1}{\lambda^2}$ & $\dfrac{\lambda}{\lambda-t}$, $t<\lambda$ \\[7pt]
$\Gamma(\alpha,\beta)$ & $\dfrac{\beta^\alpha}{\Gamma(\alpha)} x^{\alpha-1} e^{-\beta x}$, $x>0$, $\alpha,\beta>0$ & $\dfrac{\alpha}{\beta}$ & $\dfrac{\alpha}{\beta^2}$ & $\left(\dfrac{\beta}{\beta-t}\right)^{\!\alpha}$ \\[7pt]
$N(\mu,\sigma^2)$ & $\dfrac{1}{\sigma\sqrt{2\pi}} e^{-\frac{(x-\mu)^2}{2\sigma^2}}$, $\sigma>0$ & $\mu$ & $\sigma^2$ & $e^{\mu t + \sigma^2 t^2/2}$ \\[7pt]
$\chi^2(n) = \Gamma\!\left(\frac{n}{2},\frac12\right)$ & $\dfrac{x^{n/2-1}e^{-x/2}}{2^{n/2}\Gamma(n/2)}$, $x>0$, $n>0$ & $n$ & $2n$ & $(1-2t)^{-n/2}$, $t<\tfrac12$ \\[7pt]
$t(n)$ & $\dfrac{\Gamma(\frac{n+1}{2})}{\sqrt{n\pi}\,\Gamma(\frac{n}{2})}\left(1+\frac{x^2}{n}\right)^{-\frac{n+1}{2}}$, $n>0$ & $0$, $n>1$ & $\dfrac{n}{n-2}$, $n>2$ & няма MGF в отворена околност на $0$ \\
\bottomrule
\end{tabular}
\end{center}

\noindent
В таблицата $\beta$ е \emph{скорост} (rate/intensity), не scale. За $U(a,b)$
формулата за $M_X(t)$ при $t=0$ се разбира по непрекъснатост като $1$;
за $\Gamma$ е валидна при $t<\beta$, за $\chi^2$ при $t<1/2$, а за
нормалното — за всяко реално $t$. За $t$-разпределението $M_X(0)=1$ по
дефиниция, но няма нетривиална околност на нулата, в която MGF да е крайна.

\begin{itemize}\itemsep2pt
  \item $\Gamma(\alpha+1) = \alpha\,\Gamma(\alpha)$;\quad $\Gamma(n) = (n-1)!$;\quad
        $\Gamma(\tfrac12) = \sqrt{\pi}$;\quad $\Exponential(\lambda) = \Gamma(1,\lambda)$.
  \item За $X\sim\Exponential(\lambda)$, $\lambda>0$,
        $F_X(x)=0$ при $x\le0$ и $F_X(x)=1-e^{-\lambda x}$ при $x>0$;
        съответно $\mathbb P(X>x)=1$ при $x\le0$ и $e^{-\lambda x}$ при $x>0$.
        Липсата на паметта $\mathbb{P}(X > s+t \given X > s) = \mathbb{P}(X > t)$
        е за $s,t\ge0$.
  \item \emph{Стандартизиране:} при $0<\sigma^2<\infty$,
        $X \sim N(\mu,\sigma^2) \implies \dfrac{X-\mu}{\sigma} \sim N(0,1)$,
        откъдето $\mathbb{P}(X < x) = \Phi\!\left(\frac{x-\mu}{\sigma}\right)$ и
        $\Phi(-z) = 1-\Phi(z)$.
  \item \emph{Устойчивост:} сума от независими нормални е нормална
        ($\mu$ и $\sigma^2$ се събират); сума от независими $\Poi$ е $\Poi$;
        сума от независими $\Gamma$ с обща скорост $\beta>0$ е $\Gamma$
        (събира се $\alpha$).
  \item \emph{Връзки:} $Z \sim N(0,1) \implies Z^2 \sim \chi^2(1)$;\quad
        сума от $n$ независими $\chi^2(1)$ е $\chi^2(n)$;\quad
        $\dfrac{Z}{\sqrt{U/n}} \sim t(n)$ при $Z \perp\!\!\!\perp U \sim \chi^2(n)$.
\end{itemize}

\section{Числови характеристики}

\begin{itemize}\itemsep2pt
  \item $\E X = \sum_j x_j p_j$ (дискретно), $\E X = \int_{-\infty}^{\infty} x f_X(x)\,dx$ (непрекъснато),
        когато съответното очакване съществува като крайно число.
  \item \emph{Формула на статистика:} ако $\E|g(X)|<\infty$, тогава
        $\E g(X) = \sum_j g(x_j)p_j$, съответно $\int g(x) f_X(x)\,dx$
        — не е нужно да се търси разпределението на $g(X)$.
  \item \emph{Линейност:} при крайни $\E|X|$ и $\E|Y|$,
        $\E(aX+bY) = a\E X + b\E Y$; независимост не се изисква.
  \item При $\E X^2<\infty$, $\Var X = \E X^2 - (\E X)^2$ и
        $\Var(aX+b) = a^2 \Var X$.
  \item При крайни втори моменти, $\Cov(X,Y) = \E[XY] - \E X\,\E Y$;\quad
        $\Var\!\left(\sum_j X_j\right) = \sum_j \Var X_j + 2\sum_{i<j} \Cov(X_i,X_j)$.
  \item Ако $X,Y$ са независими и $\E|XY|<\infty$, тогава
        $\E[XY] = \E X\,\E Y$ и $\Cov(X,Y)=0$. Обратното \emph{не} е вярно.
  \item При $0<\Var X,\Var Y<\infty$,
        $\rho(X,Y) = \dfrac{\Cov(X,Y)}{\sqrt{\Var X}\sqrt{\Var Y}} \in [-1,1]$;\quad
        $|\rho| = 1 \iff Y = aX+b$ п.с.
  \item \emph{Очакване чрез опашките} (за $X \in \{0,1,2,\dots\}$):
        $\displaystyle \E X = \sum_{n \ge 0} \mathbb{P}(X > n)$, като равенството
        е в $[0,\infty]$; при конвенция само с крайни очаквания се изисква $\E X<\infty$.
  \item \emph{Индикатори:} $\E \ind_A = \mathbb{P}(A)$. Разбиването на брояч на сума от
        индикатори е най-прекият път към $\E X$, дори когато събираемите не са независими.
  \item \emph{Част срещу цяло:} при съществуващи втори моменти,
        $\Cov(X, X+Y) = \Var X + \Cov(X,Y)$; при $X \perp\!\!\!\perp Y$ и
        $0<\Var X,\Var Y<\infty$ —
        $\rho(X, X+Y) = \dfrac{\sigma_X}{\sqrt{\sigma_X^2+\sigma_Y^2}}$, а за независими еднакво
        разпределени $\rho(X_1, S_n) = \dfrac{1}{\sqrt n}$ (при $n=2$: $1/\sqrt2$, независимо от
        разпределението).
  \item \emph{Събираемо при фиксирана сума:} за интегрируеми независими еднакво разпределени
        $X_j$ и $S_n=\sum_jX_j$, $\E[X_1\given S_n]=S_n/n$ п.с. (за $n\ge1$).
        Записът $\E[X_1\given S_n=t]=t/n$ означава избрана регулярна условна версия
        и е само почти навсякъде по закона на $S_n$, особено когато $\mathbb P(S_n=t)=0$.
  \item \emph{Условно очакване:} при интегруем $X$,
        $\E[\E[X \given Y]] = \E X$ (формула за пълното очакване); при $\E X^2<\infty$,
        $\Var X = \E[\Var(X \given Y)] + \Var(\E[X \given Y])$.
        Вж. \crosslecture{thm:total-variance}{дефиницията и извеждането в лекция 8}.
  \item В дискретния случай $\E[X\given Y=y]$ се дефинира само за атоми
        с $\mathbb P(Y=y)>0$ чрез деление на тази маса; еднаквите стойности
        на условното средно се групират според pushforward-разпределението на $Y$.
\end{itemize}

\section{Функция на разпределение и трансформации}

\begin{itemize}\itemsep2pt
  \item В тези записки $F_X(x) = \mathbb{P}(X < x)$ (строго неравенство); $F$ е ненамаляваща,
        непрекъсната отляво, $F(-\infty)=0$, $F(+\infty)=1$.
  \item $\mathbb{P}(a \le X < b) = F_X(b) - F_X(a)$;\quad
        за непрекъсната $X$: $F_X'=f_X$ почти навсякъде (точково в точките на
        непрекъснатост на $f_X$) и $\mathbb{P}(X=a)=0$.
  \item \emph{Смяна на променливите} (едно-към-едно $C^1$-монотонна $g$ с
        $g'\ne0$ върху интервала на поддръжка $I$): ако $X$ има плътност $f_X$
        върху $I$, тогава почти навсякъде за $y\in g(I)$
        \[ f_Y(y) = f_X\big(g^{-1}(y)\big)\,\Big|\tfrac{d}{dy}g^{-1}(y)\Big| , \]
        а $f_Y(y)=0$ извън $g(I)$. По-общо е достатъчно да се изиска абсолютно
        непрекъсната обратна функция; само строгата монотонност не е достатъчна.
  \item \emph{Метод на функцията на разпределение:} намираме $F_Y(y)=\mathbb{P}(g(X)<y)$ и
        диференцираме — работи и когато $g$ не е монотонна.
  \item \emph{Двумерни:} $f_X(x) = \int f_{X,Y}(x,y)\,dy$ (маргинална);\quad
        $X \perp\!\!\!\perp Y \iff f_{X,Y}(x,y) = f_X(x) f_Y(y)$ почти навсякъде.
  \item \emph{Сума от независими:} $f_{X+Y}(z) = \int f_X(x) f_Y(z-x)\,dx$ (свъртка).
  \item \emph{Условна плътност:} при $f_X(x)>0$ (и за избрана регулярна версия),
        $f_{Y \given X}(y \given x) = f_{X,Y}(x,y)/f_X(x)$, откъдето
        $\E[Y \given X = x] = \int y\, f_{Y \given X}(y \given x)\,dy$. Практически: фиксираме
        $x$, взимаме сечението на съвместната плътност и го нормираме — знаменателят е
        точно интегралът на сечението.
  \item \emph{Смяна на променливите в две измерения.} За отворени измерими области
        и $C^1$-дифеоморфизъм между тях с $C^1$ обратно
        $(x,y)=h(z,w)$ и ненулев Якобиан:
        \[ f_{Z,W}(z,w) = f_{X,Y}\big(h(z,w)\big)\,|J|, \qquad
           J = \det\begin{pmatrix}
             \partial x/\partial z & \partial x/\partial w \\
             \partial y/\partial z & \partial y/\partial w
           \end{pmatrix}, \]
        формулата е почти навсякъде върху образа на областта и плътността е нула
        извън този образ; след което $f_Z(z) = \int f_{Z,W}(z,w)\,dw$.
        Готови якобиани:
        \[ \begin{array}{@{}lll@{}}
             Z = X+Y, & W = X: & |J| = 1; \\
             Z = XY,  & W = X: & |J| = 1/|w|; \\
             Z = X/Y, & W = Y: & |J| = |w| .
           \end{array} \]
        Границите са цялата трудност: \emph{всяко} неравенство от дефиниционната област се
        преписва в новите променливи и получената система се решава за $w$ при фиксирано $z$
        — оттам излизат и границите на интеграла, и областта на $z$. Решен пример с
        нетривиална област — \crosslecture{supp:joint-worked}{допълнението в лекция 9}.
\end{itemize}

\subsection*{Максимум и минимум на независими величини}

Нека $X_1,\dots,X_n$ са независими с функции на разпределение $F_1,\dots,F_n$, и нека
\[ M = \max_j X_j, \qquad m = \min_j X_j . \]
И двете се получават \emph{без} интегриране — само чрез преминаване към сечение на събития:
\[
  \{M < x\} = \bigcap_j \{X_j < x\}
  \;\Longrightarrow\;
  F_M(x) = \prod_{j=1}^{n} F_j(x)
  \;\overset{\text{е.р.}}{=}\; \big(F(x)\big)^{n},
\]
\[
  \{m \ge x\} = \bigcap_j \{X_j \ge x\}
  \;\Longrightarrow\;
  F_m(x) = 1 - \prod_{j=1}^{n}\big(1 - F_j(x)\big)
  \;\overset{\text{е.р.}}{=}\; 1 - \big(1-F(x)\big)^{n}.
\]
Плътностите се намират чрез диференциране; при еднакво разпределени
\[
  f_M(x) = n\,\big(F(x)\big)^{n-1} f(x), \qquad
  f_m(x) = n\,\big(1-F(x)\big)^{n-1} f(x).
\]

\begin{itemize}\itemsep2pt
  \item Работи се с $M$ през $F$, а с $m$ през \emph{опашката} $1-F$ — това е цялата идея.
        За минимум директното разписване на обединение води до включване-изключване и е
        излишно трудно.
  \item \emph{Полезен частен случай:} за краен брой независими
        $\Exponential(\lambda_j)$ с $\lambda_j>0$ минимумът е
        отново експоненциален, $m \sim \Exponential\!\left(\sum_j \lambda_j\right)$; при
        еднакви $\lambda$ това е $\Exponential(n\lambda)$. Максимумът \emph{не} е
        експоненциален.
  \item \emph{Пример:} за $n\ge1$, $\theta>0$ и $X_j \sim U(0,\theta)$ имаме
        $F_M(x) = (x/\theta)^n$ и
        $f_M(x) = n x^{n-1}/\theta^n$ — точно оценката $X^*$ от метода на максималното
        правдоподобие в лекция~12.
  \item \emph{Обща поредна статистика:} $k$-тата най-малка стойност сред $n$ е.р.\ наблюдения има
        плътност
        $\dfrac{n!}{(k-1)!\,(n-k)!}\,\big(F(x)\big)^{k-1}\big(1-F(x)\big)^{n-k} f(x)$.
        За $U(0,1)$ това е $\text{Beta}(k,\,n-k+1)$.
\end{itemize}

\section{Пораждащи функции и функции на моментите}

\begin{itemize}\itemsep2pt
  \item $g_X(s) = \E s^X$ (за целочислени $X \ge 0$);\quad $M_X(t) = \E e^{tX}$,
        като за MGF се посочва областта от $t$, където очакването е крайно.
  \item В разширен смисъл $\E X=g_X'(1-)$; ако $\E X^2<\infty$, тогава
        $\Var X = g_X''(1-) + g_X'(1-) - (g_X'(1-))^2$;\quad
        при съществуваща $k$-та производна в нулата
        $k!\,\mathbb{P}(X=k) = g_X^{(k)}(0)$.
  \item Ако MGF е крайна в отворена околност на $0$, тогава
        $\E X^k = M_X^{(k)}(0)$ и равенството на MGF-ите в такава околност
        характеризира разпределението. За PGF равенството на функциите в
        общата област на сходимост също характеризира разпределението.
  \item \emph{Суми от независими:} $g_{X+Y} = g_X \cdot g_Y$ и
        $M_{X+Y} = M_X \cdot M_Y$ във всяка обща област, където произведенията
        са крайни.
        Това е основният начин да се разпознае разпределението на сума.
\end{itemize}

\section{Неравенства и гранични теореми}

\begin{itemize}\itemsep2pt
  \item \emph{Марков:} при $X \ge 0$, $a>0$ и $\E X<\infty$,
        $\mathbb{P}(X > a) \le \dfrac{\E X}{a}$.
  \item \emph{Чебишов:} при $\E X^2<\infty$ и $a>0$,
        $\mathbb{P}(|X - \E X| > a) \le \dfrac{\Var X}{a^2}$.
  \item \emph{Слаб ЗГЧ:} за iid $X_j$ с $\E|X_1|<\infty$,
        $\dfrac{S_n}{n} \xrightarrow{\mathbb{P}} \mu$;\quad
        \emph{усилен ЗГЧ:} при същите предпоставки
        $\dfrac{S_n}{n} \xrightarrow{\text{п.с.}} \mu$.
  \item \emph{ЦГТ:} за iid $X_j$ с $\mu=\E X_1$ и $0<\Var X_1=\sigma^2<\infty$,
        $\displaystyle \frac{S_n - n\mu}{\sigma\sqrt{n}} \xrightarrow{d} N(0,1)$, откъдето
        $\mathbb{P}(S_n < b) \approx \Phi\!\left(\dfrac{b-n\mu}{\sigma\sqrt{n}}\right)$.
        Показаната MGF-аргументация е само доказателство-скица при MGF в
        околност на нулата; общата ЦГТ се доказва например с характеристични функции.
  \item \emph{Йерархия:} п.с. $\implies$ по вероятност $\implies$ по разпределение;
        обратните посоки не са верни.
\end{itemize}

\section{Статистика}

\begin{itemize}\itemsep2pt
  \item \emph{Извадкови характеристики:} за $n\ge1$,
        $\overline{X}_n = \frac1n\sum X_j$; за $n\ge2$,
        $S^2 = \frac{1}{n-1}\sum (X_j - \overline{X}_n)^2$ (неизместена);\quad
        $\hat\sigma^2 = \frac{n-1}{n}S^2$ (МПО, изместена).
  \item При iid извадка с $\E X_1=\mu$ и $\Var X_1=\sigma^2<\infty$,
        $\E \overline{X}_n = \mu$, $\Var \overline{X}_n = \dfrac{\sigma^2}{n}$;\quad
        при $n\ge2$ и $\E X_1^2<\infty$, $\E S^2 = \sigma^2$.
  \item \emph{Максимално правдоподобие:} за фиксирани данни $\vec x$ максимизираме
        $L(\theta;\vec x)=\prod_j g_\theta(x_j)$, където $g_\theta$ е pmf или плътност,
        обикновено през $\ln L$ в областта $L>0$. Уравнението на score дава само
        вътрешни кандидати; проверяват се границата и достигането на супремума.
        Ако носителят зависи от $\theta$ (напр.\ $U(0,\theta)$), решава цялата форма на $L$.
  \item \emph{Метод на моментите:} приравняваме теоретични към извадкови моменти.
  \item \emph{Централна статистика} $T(\vec X,\theta)$: с разпределение,
        което не зависи от $\theta$, и строго монотонна по $\theta$ върху параметърното
        пространство, така че уравненията $q_i=T(\vec X,\theta)$ имат решения в
        допустимия диапазон. Оттук интервалът се получава чрез обръщане на
        $\mathbb{P}(q_1 < T < q_2) = \gamma$; при липса на обратимост се използва
        обобщена обратна функция и се проверява покритието отделно.
        Получените граници са доверителни само ако
        $\mathbb P_\theta(I_1(\vec X)<\theta<I_2(\vec X))\ge\gamma$ за всяко
        $\theta$ от параметърното пространство.
  \item При iid $N(\mu,\sigma^2)$ наблюдения, $n\ge1$ и известна $\sigma>0$:\;
        $\overline{X}_n \pm q_{\frac{1+\gamma}{2}}\dfrac{\sigma}{\sqrt{n}}$.
  \item При iid $N(\mu,\sigma^2)$ наблюдения, $n\ge2$, $\sigma>0$ и \emph{неизвестна} $\sigma$:\;
        $\overline{X}_n \pm t_{\frac{1+\gamma}{2},\,n-1}\dfrac{S}{\sqrt{n}}$, понеже
        $\dfrac{\overline{X}_n-\mu}{S/\sqrt{n}} \sim t(n-1)$.
  \item При iid $N(\mu,\sigma^2)$, $n\ge1$, $\sigma>0$ и известна $\mu$,
        $\dfrac{\sum_{j=1}^n(X_j-\mu)^2}{\sigma^2}\sim\chi^2(n)$.
        При неизвестна $\mu$, $n\ge2$,
        $\dfrac{(n-1)S^2}{\sigma^2} \sim \chi^2(n-1)$ — интервалът използва
        \emph{два различни} квантила, защото $\chi^2$ не е симетрично.
  \item \emph{Проверка на хипотези:} грешка от I род $\alpha = \mathbb{P}(\vec X \in W \given H_0)$,
        от II род $\beta = \mathbb{P}(\vec X \notin W \given H_1)$. Лемата на
        Нейман\,--\,Пиърсън дава най-мощната област чрез отношението
        $L_1(\vec x) \ge k\,L_0(\vec x)$.
  \item \emph{Линейна регресия:} при фиксиран дизайн, независими
        $N(0,\sigma^2)$ грешки с $\sigma^2>0$, $n\ge2$ и
        $A=\sum (X_k-\overline X)^2>0$,
        $b_1 = \dfrac{\sum (X_k-\overline X)(Y_k-\overline Y)}{A}$,\;
        $b_0 = \overline Y - b_1 \overline X$. За остатъчната дисперсия и $t$
        изводи е нужно $n\ge3$ (остатъчни степени $n-2>0$); тук
        $\overline X,\overline Y$ са аритметични извадкови средни.
\end{itemize}

\section{Апарат от анализа}

Изброеното по-долу не е ново — то е от курса по анализ — но точно тези няколко неща се
появяват отново и отново в задачите по вероятности. До всяко е отбелязано къде се използва.

\subsection*{Редове}

\begin{itemize}\itemsep3pt
  \item \emph{Геометричен ред и двете му производни} (за $|q|<1$) — работният кон при
        $\Geo$ и $\NegBin$:
        \[
          \sum_{k \ge 0} q^k = \frac{1}{1-q}, \qquad
          \sum_{k \ge 1} k q^{k-1} = \frac{1}{(1-q)^2}, \qquad
          \sum_{k \ge 2} k(k-1) q^{k-2} = \frac{2}{(1-q)^3} .
        \]
        Вторите две се получават от първата чрез \emph{почленно диференциране} — законно е
        във вътрешността на интервала на сходимост. Оттук идват $\E X$ и $\Var X$ през
        $g_X'(1)$ и $g_X''(1)$.
  \item \emph{Експоненциален ред:} $\displaystyle \sum_{k \ge 0} \frac{x^k}{k!} = e^{x}$.
        С него се проверява, че поасоновите вероятности се сумират до $1$, и се получава
        $g_X(s) = e^{\lambda(s-1)}$.
  \item \emph{Отрицателен биномен ред:}
        $\displaystyle \sum_{k \ge 0} \binom{r+k-1}{k} q^k = \frac{1}{(1-q)^r}$ —
        точно това превръща $g_X(s) = \left(\frac{p}{1-qs}\right)^r$ обратно във
        вероятностите на $\NegBin(r,p)$.
  \item \emph{Бином на Нютон:} $(a+b)^n = \sum_k \binom{n}{k}a^k b^{n-k}$ — оттук е
        $g_X(s) = (q+ps)^n$ за биномното.
  \item $\displaystyle \sum_{n \ge 1} \frac{1}{n^2} = \frac{\pi^2}{6}$ — появява се при
        конструиране на разпределения с тежки опашки (нормиращата константа е $6/\pi^2$).
  \item \emph{Смяна на реда на сумиране} в двойна сума: законна при неотрицателни
        събираеми. Това е целият трик зад $\E X = \sum_{n \ge 0} \mathbb{P}(X>n)$.
\end{itemize}

\subsection*{Граници}

\begin{itemize}\itemsep3pt
  \item $\displaystyle \lim_{n \to \infty}\Big(1 + \frac{x}{n}\Big)^{n} = e^{x}$,\;
        в частност $\big(1-\frac{\lambda}{n}\big)^n \to e^{-\lambda}$ — това е моторът на
        теоремата на Поасон ($\Bin(n,\lambda/n) \to \Poi(\lambda)$).
  \item \emph{Формула на Стирлинг:} $n! \sim \sqrt{2\pi n}\,\big(\frac{n}{e}\big)^{n}$ —
        за асимптотика на комбинаторни изрази и на централни биномни коефициенти.
  \item При неопределеност от вида $\frac{0}{0}$ или $\frac{\infty}{\infty}$ —
        правилото на Лопитал; най-често се налага при проверка на гранични разпределения.
\end{itemize}

\subsection*{Интеграли}

\begin{itemize}\itemsep3pt
  \item \emph{Смяна на променливите} е основният инструмент. Ако $u = g(x)$, то
        $\int f(g(x))g'(x)\,dx = \int f(u)\,du$; при определен интеграл се сменят и
        границите. Типични случаи: $u = \frac{x-\mu}{\sigma}$ (стандартизиране),
        $u = \beta x$ или $u = x(\beta-t)$ (свеждане към гама-интеграл),
        $u = x^2/2$ (нормални опашки).
  \item \emph{Гама-интеграл:}
        $\displaystyle \int_0^{\infty} x^{\alpha-1} e^{-\beta x}\,dx = \frac{\Gamma(\alpha)}{\beta^{\alpha}}$.
        Почти всяко пресмятане с $\Gamma$, $\Exponential$ или $\chi^2$ се свежда до него.
        Свойства: $\Gamma(\alpha+1) = \alpha\Gamma(\alpha)$, $\Gamma(n) = (n-1)!$,
        $\Gamma(\tfrac12) = \sqrt{\pi}$.
  \item \emph{Гаусов интеграл:}
        $\displaystyle \int_{-\infty}^{\infty} e^{-x^2/2}\,dx = \sqrt{2\pi}$ — оттук е
        нормиращата константа на нормалното разпределение.
  \item \emph{Интегриране по части:} $\int u\,dv = uv - \int v\,du$. Използва се при
        моменти на непрекъснати величини и за да се докаже, че за $X \ge 0$
        \[ \E X = \int_0^{\infty} \big(1 - F_X(x)\big)\,dx \]
        — непрекъснатият аналог на сумата по опашките.
  \item \emph{Диференциране под знака на интеграла} (по параметър): при достатъчно добри
        условия $\frac{d}{d\theta}\int f(x,\theta)dx = \int \frac{\partial f}{\partial \theta}dx$.
        С този похват се пресмята $\E Z^2$ за $N(0,1)$, без да се интегрира по части.
\end{itemize}

\subsection*{Похвати, които се повтарят в задачите}

\begin{itemize}\itemsep3pt
  \item \emph{Намиране на нормираща константа:} условието $\int f = 1$ (или $\sum p_k = 1$)
        е уравнение за неизвестния множител — почти всяка задача с „намерете $c$“ е това.
  \item \emph{Моменти:} $\E X = \int x f$, $\E X^2 = \int x^2 f$, после
        $\Var X = \E X^2 - (\E X)^2$. Търсенето на $\Var$ директно по дефиниция е по-дълго.
  \item \emph{Логаритмуване преди диференциране:} при максимално правдоподобие се
        максимизира $\ln L$, защото произведението става сума. Максимумът е същият, понеже
        $\ln$ е строго растяща.
  \item \emph{Проверка на плътност:} $f \ge 0$ и $\int f = 1$ — двете заедно са
        достатъчни, за да е $f$ плътност на някаква величина.
  \item \emph{Симетрия:} ако подинтегралната функция е нечетна и границите са симетрични,
        интегралът е нула. Така се вижда веднага, че $\E Z = 0$ за $N(0,1)$.
\end{itemize}
